\documentclass[11pt]{article}

% basic packages
\usepackage[margin=1in]{geometry}
\usepackage[pdftex]{graphicx}
\usepackage{amsmath,amssymb,amsthm}
\usepackage{william}

% page formatting
\usepackage{fancyhdr}
\pagestyle{fancy}

\renewcommand{\sectionmark}[1]{\markright{\textsf{\arabic{section}. #1}}}
\renewcommand{\subsectionmark}[1]{}
\lhead{\textbf{\thepage} \ \ \nouppercase{\rightmark}}
\chead{}
\rhead{}
\lfoot{}
\cfoot{}
\rfoot{}
\setlength{\headheight}{14pt}

\linespread{1.03}
\setlength{\parindent}{0.2in}
\setlength{\emergencystretch}{3em} % last-resort stretch: avoid overfull boxes from long inline math
\setcounter{secnumdepth}{2}
\setcounter{tocdepth}{2}

% ---- local macros for this explainer ----
\DeclareMathOperator{\Spec}{Spec}
\DeclareMathOperator{\Ann}{Ann}
\DeclareMathOperator{\Frac}{Frac}
\DeclareMathOperator{\Nil}{Nil}
\DeclareMathOperator{\Jac}{Jac}
\DeclareMathOperator{\Supp}{Supp}
\DeclareMathOperator{\Ass}{Ass}
\DeclareMathOperator{\Tor}{Tor}
\DeclareMathOperator{\Ext}{Ext}
\DeclareMathOperator{\gr}{gr}
\newcommand{\mfp}{\mathfrak{p}}
\newcommand{\mfm}{\mathfrak{m}}
\newcommand{\mfq}{\mathfrak{q}}
\newcommand{\mfa}{\mathfrak{a}}
\newcommand{\mfn}{\mathfrak{n}}
\newcommand{\dirlim}{\varinjlim}
\newcommand{\invlim}{\varprojlim}
% green "Exercise" box, reusing the green-box style from william.sty
\declaretheorem[style=thmgreenbox,name=Exercise,numberwithin=section]{xca}

% a small helper for the per-section reading pointers
\newcommand{\sources}[1]{\smallskip\noin{\small\emph{Read alongside:} #1}\par\medskip}

\begin{document}

\thispagestyle{empty}
\bigskip
\vspace{0.1cm}

\begin{center}
{\fontsize{20}{20} \selectfont An Explainer on}
\vskip 14pt
{\fontsize{34}{34} \selectfont \bf \sffamily Commutative Algebra}
\vskip 10pt
{\fontsize{16}{16}\selectfont \sffamily A learning path from the modern, categorical point of view}
\vskip 22pt
{\fontsize{16}{16} \selectfont \rmfamily \href{https://will-lancer.github.io}{Will Lancer}}
\vskip 6pt
{\fontsize{12}{12} \selectfont \ttfamily will.m.lancer@gmail.com}
\vskip 18pt
\end{center}

{\parindent0pt \baselineskip=15.5pt}
\noin
Commutative algebra is two subjects wearing one coat. It is the \emph{local algebra of
geometry}---the thing that tells you what a variety looks like near a point---and it is
the \emph{arithmetic of numbers}---the thing that tells you how primes factor in a ring of
integers. The modern, Grothendieck-era insight is that these are the \emph{same} subject,
because a commutative ring \emph{is} a space: its points are prime ideals, its functions
are its elements, and a map of rings is a map of spaces running the other way. These notes
are a guided path to that point of view. We learn the subject the way it is actually
organized today: \vocab{$\Spec$} (the spectrum) as the central object, \vocab{modules} and
\vocab{localization} as the two everyday tools, and the language of \vocab{categories,
universal properties, and adjoint functors} as the connective tissue that makes the whole
thing hang together.

\medskip
\noin
\emph{What this document is.} It is the \textbf{recommended learning path}---an annotated
table of contents. Each section says \emph{what} it covers, \emph{why} it sits where it
does in the modern arc, and \emph{what to read} for it in the three primary sources (see
the reference list). Read it as a syllabus and a map; the full worked notes are built out
section by section on top of this skeleton.

\medskip
\noin
\emph{Prerequisites.} A first course in abstract algebra (groups, rings, fields), linear
algebra, and a little point-set topology. \emph{No} category theory is assumed---we build
it from scratch in \S2, exactly where it is first needed.

\medskip
\noin
\emph{How to read.} Boxes labelled \textsf{Definition}/\textsf{Theorem} are the
load-bearing statements; \textsf{Intuition} and \textsf{Note} boxes are the ``how a working
algebraist pictures it'' commentary; yellow \textsf{Equation} boxes hold a ``carry this in
your head'' slogan. Green \textsf{Exercise} boxes are scattered throughout, tagged
\textbf{(easy)}, \textbf{(classic)}, \textbf{(medium)}, or \textbf{(thinker)}---do them as
you go. An end-of-path exercise set and an annotated ``books to own'' list close the notes.

\newpage
\microtoc
\newpage

% ======================================================================
\section{The big picture: commutative algebra as geometry and arithmetic}
% ======================================================================

\noindent
Before a single definition, we answer the only question that matters at the start: why does
this subject exist, and what is it secretly about? The answer is a dictionary. To every
commutative ring $R$ we attach a space $\Spec R$ whose points are the prime ideals of $R$;
the elements of $R$ become functions on that space, and a homomorphism $R\to S$ becomes a
continuous map $\Spec S \to \Spec R$. Geometry (varieties, their singularities, families)
and arithmetic (factorization, ramification, class groups) turn out to be two readings of
this one dictionary. This opening section sets up the dictionary, names the three tools we
will lean on the whole way through, and lays out the route.

\subsection{Two faces of one subject: the local algebra of geometry, the arithmetic of numbers}

Start with two questions you can already feel the weight of, one from geometry and one from
arithmetic.

\emph{The geometric question.} Take a plane curve, say
$C\colon y^2 = x^3 - x$. How do you study it \emph{algebraically}? The natural move is to
look at the polynomial functions on the plane and see which two agree on $C$: two
polynomials restrict to the same function on $C$ exactly when their difference is a multiple
of $y^2 - x^3 + x$. So the functions on $C$ form the quotient ring
$k[x,y]/(y^2-x^3+x)$. Every geometric feature of $C$---where it is smooth, where it crosses
itself, how many pieces it has---turns out to be readable off this ring. Geometry has become
algebra.

\emph{The arithmetic question.} In $\ints$ every integer factors uniquely into primes. Now
enlarge $\ints$ to $\ints[\sqrt{-5}] = \{a + b\sqrt{-5} : a,b\in\ints\}$. There,
\[
6 = 2\cdot 3 = (1+\sqrt{-5})(1-\sqrt{-5}),
\]
and all four factors are irreducible: unique factorization \emph{fails}. Dedekind's
resolution---one of the births of the subject---was to stop factoring \emph{elements} and
start factoring \emph{ideals}, in which language $(6)$ does factor uniquely, and to measure
the failure by an invariant of the ring (the class group). Arithmetic has become algebra
about ideals.

The point is that both questions are about a \vocab{commutative ring}: in the first it is a
ring of \emph{functions on a space}, in the second a ring of \emph{numbers}. Commutative
algebra is the claim that these are not two subjects but one.

\begin{iidea}[The thesis of commutative algebra]
A commutative ring is, simultaneously, the ``local algebra of geometry'' (the functions on a
space) and the ``arithmetic of a number system.'' The two readings are unified because both
are commutative rings, and the \emph{same} tools---ideals, modules, localization---resolve
both the geometric and the arithmetic questions. Learning to read a ring in both dialects at
once is the whole skill.
\end{iidea}

The unification is not vague analogy; it is a precise dictionary, and the cleanest place to
see it is the parallel between the integers $\ints$ and the polynomial ring $k[x]$ in one
variable over a field $k$.\footnote{This $\ints \leftrightarrow k[x]$ parallel is the
backbone of the whole modern viewpoint; A--M open their book by announcing that the
``central notion in commutative algebra is that of a prime ideal'' and that, following the
modern trend, they emphasize modules and localization (A--M, Introduction). We take that
emphasis as our organizing principle.} Both are \vocab{principal ideal domains}: every ideal
is generated by one element. In $\ints$ the ``irreducible'' objects are the prime numbers
$p$; in $k[x]$ (take $k$ algebraically closed for a clean statement) they are the linear
polynomials $x-a$, one for each point $a$ of the line. Read the table across, and the
arithmetic column and the geometry column are the same story told twice.

\begin{center}
\begin{tabular}{lll}
\hline
 & \textbf{Arithmetic: } $\ints$ & \textbf{Geometry: } $k[x]$ \\
\hline
the ring & integers $\ints$ & polynomials $k[x]$ \\
an irreducible & prime number $p$ & linear factor $x-a$ \\
a ``point'' & prime ideal $(p)$ & prime ideal $(x-a)$ \\
residue field there & $\ints/(p)=\mathbb{F}_p$ & $k[x]/(x-a)=k$ \\
``value'' of an element & $n \bmod p$ & $f(a)$ \\
the generic point & $(0)$, with field $\mathbb{Q}$ & $(0)$, with field $k(x)$ \\
\hline
\end{tabular}
\end{center}

\noindent
The last two rows are the punchline, and they are worth pausing on. Reducing an integer $n$
modulo a prime $p$ and evaluating a polynomial $f$ at a point $a$ \emph{are the same
operation}: in each case you take an element of the ring and pass to the quotient by the
``point'' to read off its value there. An integer is, in this precise sense, a function on
the prime numbers, whose value at $p$ is $n \bmod p$. That single realization---numbers are
functions, primes are points---is the seed from which the entire modern subject grows. The
next subsection makes it a definition.

A good analogy earns its keep by showing you exactly where it frays, so before we lean on
this one we should say where it does.

\begin{reemark}[Where the analogy frays---and the deeper analogy underneath]
The $\ints\leftrightarrow k[x]$ dictionary is a guide, not an identity, and its seams are
themselves famous. Two honest disanalogies. \emph{First}, the residue fields of $k[x]$ are
all the same field $k$, whereas those of $\ints$ are the \emph{varying} finite fields
$\mathbb{F}_p$---so $k[x]$ is really the \emph{geometric} (``equal characteristic'') analogue
of $\ints$. The closer \emph{arithmetic} analogue is the polynomial ring $\mathbb{F}_q[t]$
over a finite field, whose residue fields $\mathbb{F}_q[t]/(\pi)$ are again finite. \emph{Second},
$\ints$ secretly carries one extra ``point'': the ordinary absolute value on $\mathbb{Q}$ is
a further, \vocab{archimedean} place with no algebraic counterpart, so $\Spec\ints$ behaves
like an \emph{affine} curve that is missing its point at infinity---supplying it is the
starting gun of Arakelov geometry. Far from being a defect, the full analogy
\[
\{\text{number fields}\}\ \longleftrightarrow\ \{\text{function fields of curves over a finite field } \mathbb{F}_q\},
\]
with a ring of integers $\mathcal{O}_K$ playing the role of (the affine part of) a curve, is
the backbone of modern number theory: it is the setting of Weil's proof of the Riemann
hypothesis for curves and of the function-field Langlands program of Drinfeld and L.\
Lafforgue. We will keep meeting it.
\end{reemark}

\subsection{The central dogma: rings are spaces, primes are points, functions are elements}

Here is the dictionary in full. The precise definitions and proofs come in \S3; right now we
want the shape of the idea so that everything afterward has a place to land.

\begin{ppreview}[The ring--space dictionary (to be made precise in \S3)]
To a commutative ring $R$ we attach a space, the \vocab{spectrum} $\Spec R$, whose
\emph{points} are the \vocab{prime ideals} of $R$.
\begin{itemize}\setlength{\itemsep}{2pt}
\item An element $f\in R$ is a \emph{function} on $\Spec R$. Its \emph{value} at a point
$\mfp$ is the image of $f$ in the \vocab{residue field}
$\kappa(\mfp)\defeq \Frac(R/\mfp)$---the fraction field of the (integral domain) quotient
$R/\mfp$.\footnote{$R/\mfp$ is a domain precisely because $\mfp$ is prime; that is the whole
reason primes, not arbitrary ideals, are the points. See \S3.} In particular, $f$
\vocab{vanishes at} $\mfp$ (its value there is $0$) exactly when $f\in\mfp$, so the points
where $f$ is zero are precisely the primes that contain it---this is what will turn
``solving $f=0$'' into the theory of ideals.
\item A ring homomorphism $\varphi\colon R\to S$ induces a map of spaces in the
\emph{opposite} direction, $\Spec S\to\Spec R$, sending a prime $\mfq\subseteq S$ to its
preimage $\varphi^{-1}(\mfq)\subseteq R$.
\end{itemize}
\end{ppreview}

Two design decisions in that box deserve their motivation, because each is a place where a
beginner's first guess is wrong.

\emph{Why prime ideals, and not just maximal ones?} The classical points of a space---a
single number $a$, a single point of the plane---correspond to \emph{maximal} ideals (this
is the Nullstellensatz, \S10). So why widen ``point'' to all primes? Two reasons, one formal
and one geometric. Formally: the preimage of a prime ideal under a ring map is always
prime, but the preimage of a \emph{maximal} ideal need not be maximal.\footnote{Example:
$\ints\hookrightarrow\mathbb{Q}$. The ideal $(0)$ is maximal in the field $\mathbb{Q}$, but
its preimage $(0)\subseteq\ints$ is prime and not maximal. So ``maximal ideals'' do not form
a functorial notion of point, while primes do. This is exactly why $\Spec$ uses primes.} If
we want $\Spec$ to be a \emph{functor}---so that algebra translates to geometry
automatically---we are forced to use primes. There is also a vivid reason the \emph{prime}
condition is the right one. The defining property $fg\in\mfp\Rightarrow f\in\mfp$ or
$g\in\mfp$ says, in the language of values, that a product of functions can vanish at $\mfp$
only if one of the factors already does; the vanishing at $\mfp$ cannot be split into two
independent pieces, so $\mfp$ behaves like an \emph{irreducible} locus---and, dually, its
ring of functions $R/\mfp$ has no zero-divisors. ``Prime ideal'' is the algebra of
``irreducible piece of the space.'' Geometrically: a non-maximal prime is the
\vocab{generic point} of an irreducible subspace. In $k[x,y]$ the prime $(y-x^2)$ is not a
single point of the plane; it is the entire parabola $y=x^2$, packaged as one point of
$\Spec k[x,y]$. Allowing irreducible subvarieties to \emph{be} points is one of the most
useful moves in all of geometry, and it is free once you take primes seriously.

\emph{Why does the map of spaces go backwards?} Because functions pull back. If
$\varphi\colon R\to S$ presents $\Spec S$ as a space mapping to $\Spec R$, then a function on
the target ($f\in R$) becomes a function on the source ($\varphi(f)\in S$) by composition.
Pulling functions back is contravariant, so the spaces map the opposite way to the rings.
This is the same variance you already know from topology: a continuous map $X\to Y$ induces
$C(Y)\to C(X)$ on rings of continuous functions.

\begin{eequation}[The slogan to carry in your head]
\[
\underbrace{\text{value of } f \text{ at the point } \mfp}_{\text{image of } f \text{ in } \kappa(\mfp)}
\qquad\rightsquigarrow\qquad
\underbrace{n \bmod p}_{\text{in } \ints}
\;\;=\;\;
\underbrace{f(a)}_{\text{in } k[x]}.
\]
Elements are functions; primes are points; reduction modulo a prime is evaluation.
\end{eequation}

\begin{eexample}[The two spectra to know by heart]
\textbf{(a) $\Spec\ints$.} The primes of $\ints$ are $(0)$ and $(p)$ for each prime number
$p$. The points $(p)$ are ``closed'' (ordinary points), with residue field
$\kappa\big((p)\big)=\mathbb{F}_p$; the value of $n$ there is $n\bmod p$. The point $(0)$ is
the \emph{generic point}: its residue field is $\kappa\big((0)\big)=\mathbb{Q}$, and the
value of $n$ there is $n$ itself. So $\Spec\ints$ is the prime numbers together with one
extra ``smeared-out'' point sitting over all of them.

\textbf{(b) $\Spec k[x]$, $k=\bar k$.} The primes are $(0)$ and $(x-a)$ for $a\in k$. The
closed points $(x-a)$ \emph{are} the line $k$; the residue field is $k$ and the value of $f$
at $(x-a)$ is $f(a)$---honest evaluation. Again there is one generic point $(0)$, with
residue field $k(x)$, the field of rational functions.
\end{eexample}

\begin{iintuition}[Elements are functions, but stranger]
An element behaves like a function on $\Spec R$, with two twists that ordinary functions
lack. First, its values live in \emph{different fields at different points} (the residue
field $\kappa(\mfp)$ changes with $\mfp$), so ``$f$'' is a function whose codomain varies.
Second, an element can be nonzero yet vanish at points, and---most strikingly---a nonzero
\vocab{nilpotent} can vanish at \emph{every} point: in $k[x]/(x^2)$ the element $x$ is not
$0$, but its value in every residue field is $0$. Classical functions cannot do this; the
ring remembers an ``infinitesimal'' direction that the naive set of points forgets. We keep
nilpotents (rather than quotient them away) precisely because they record this
infinitesimal data---the ring $k[\varepsilon]/(\varepsilon^2)$ is a point carrying a tangent
vector. This is the first hint of why \emph{scheme} theory works with rings, not just point
sets.
\end{iintuition}

The contravariant map of spectra is not a technicality; it is where arithmetic turns into the
geometry of covering maps.

\begin{eexample}[{A ring map as a branched cover: $\ints\hookrightarrow\ints[i]$}]
The inclusion $\ints\hookrightarrow\ints[i]$ of the integers into the Gaussian integers
induces $\Spec\ints[i]\to\Spec\ints$. Over a point $(p)$ of $\Spec\ints$ sits the
\emph{fiber}: the primes of $\ints[i]$ lying above $p$, i.e.\ the way $p$ factors in
$\ints[i]$. And the classical trichotomy of number theory is exactly a statement about this
map:
\[
(5)=(2+i)(2-i)\ \text{\small(splits: two points)},\quad
(3)\ \text{\small(inert: one point)},\quad
(2)=-i(1+i)^2\ \text{\small(ramifies)}.
\]
So $\Spec\ints[i]\to\Spec\ints$ is a (branched) double cover\footnote{The cover is branched
over only finitely many primes: $2$ is the sole \emph{ramified} prime here, and it is exactly
the one dividing the discriminant $-4$ of $\ints[i]$. Ramification is where the two sheets of
the cover collide, and it is detected algebraically by the discriminant---the same
discriminant you met under the square root in the quadratic formula.} of $\Spec\ints$, and
``how $p$ splits'' is the same data as ``what the fiber looks like.'' The arithmetic of the
Gaussian integers \emph{is} the geometry of this covering map---the prototype of everything
in algebraic number theory.
\end{eexample}

\begin{xca}[\textbf{(easy)} preimages of primes]
Let $\varphi\colon R\to S$ be a ring homomorphism and $\mfq\subseteq S$ a prime ideal. Show
$\varphi^{-1}(\mfq)$ is a prime ideal of $R$. \emph{(Hint: $R/\varphi^{-1}(\mfq)$ injects
into $S/\mfq$.)} Then give an example showing the preimage of a \emph{maximal} ideal need
not be maximal.
\end{xca}

\subsection{The three pillars of the modern view: modules, localization, categories}

What separates the modern treatment of this subject from its classical ancestor is which
tools it treats as fundamental. Three of them carry almost all the weight, and naming them
now---with their motivation---turns the rest of these notes into a guided tour rather than a
list of theorems.

\textbf{Pillar 1: modules---linear algebra over a ring.} You understand a group by its
representations; you understand a ring by the things it acts on, its \vocab{modules}. A
module is ``a vector space whose scalars come from a ring instead of a field,'' and almost
every object we have met is one: ideals, quotient rings, and the ring itself are all
$R$-modules; abelian groups are exactly $\ints$-modules; vector spaces are exactly modules
over a field. The modern slogan is that commutative algebra is the study of the category of
$R$-modules. The single phenomenon that makes the theory deep is that, unlike vector spaces,
\emph{a module need not have a basis}---and measuring how badly bases fail (via free
resolutions, $\Tor$, $\Ext$) is a large part of the subject. A vivid case to keep in mind: a
module over $k[x]$ is a vector space equipped with a single linear operator (the action of
$x$), so the classification of $k[x]$-modules \emph{is} the theory of Jordan and rational
canonical forms.\footnote{The same slogan reaches much further: a module over a group ring
$k[G]$ is precisely a representation of $G$, so representation theory is ``the module theory
of $k[G]$.'' That is one of the doorways from commutative algebra toward the Langlands
program, where representations are the central objects.} Linear algebra is already
commutative algebra, in disguise.

\textbf{Pillar 2: localization---geometry is local.} To understand a space you look near each
point; the algebraic counterpart is \vocab{localization}, the operation of inverting
elements. Inverting everything outside a prime $\mfp$ produces the \vocab{local ring}
$R_\mfp$, which sees only the behavior of $R$ near the point $\mfp$---it is the ring of
germs of functions there, exactly as in the theory of manifolds. The reward is the
\vocab{local--global principle}: a great many properties hold for $R$ if and only if they
hold at every prime, so global theorems can be proved one point at a time. Localization is
the technique that makes ``study the space near a point'' into algebra.

\textbf{Pillar 3: categories---the connective language.} We will constantly build new rings
and modules out of old (quotients, products, tensor products, localizations, completions,
free objects) and will want to say that each construction is \emph{canonical}, pin it down
\emph{up to unique isomorphism}, and track how it interacts with maps. The precise language
for all of this is category theory: \vocab{universal properties}, \vocab{functors}, and
above all \vocab{adjoint functors}. It is not decoration. The fact that the tensor product
is the left adjoint of $\Hom$, and that ``left adjoints are right exact,'' organizes
\S\S5--8 in a single stroke.\footnote{This is the explicit thesis of Altman--Kleiman, who
redo A--M with the categorical machinery in front: ``Tensor Product is the left adjoint of
Hom \dots\ any left adjoint preserves arbitrary direct limits'' (A--K, Preface). We develop
exactly this language, from scratch, in \S2.} Section 2 builds this language from nothing.

\begin{eequation}[The three pillars]
\[
\boxed{\ \textbf{modules}\ }\qquad \boxed{\ \textbf{localization}\ }\qquad \boxed{\ \textbf{categories}\ }
\]
Study a ring through what it acts on (modules); study it one point at a time (localization);
and keep track of every construction by the maps it represents (categories).
\end{eequation}

\subsection{If you remember one thing; and a map of the route ahead}

\begin{iidea}[If you remember one thing]
A commutative ring is the ring of functions on a space, its spectrum $\Spec R$, whose points
are the prime ideals. Doing algebra to the ring \emph{is} doing geometry to the space:
ideals cut out subspaces, quotients restrict to them, localization zooms in, and a ring map
is a map of spaces running the other way. Every definition that follows is a tool for making
this correspondence precise and computable---and, read in the other dialect, for turning
arithmetic into the geometry of $\Spec\ints$.
\end{iidea}

With the thesis in hand, here is the shape of the journey, so you can see where each section
fits in the arc.

\begin{ppreview}[A map of the route ahead]
We first learn the \emph{language} (\S2: categories, universal properties, adjoints), then
build the central \emph{objects and the dictionary} (\S3: $\Spec R$; \S4: modules; \S5:
exact sequences). Next come the two everyday \emph{tools}: localization (\S6) and tensor
products with flatness (\S7), with their homological shadow $\Tor$/$\Ext$ (\S8). We then do
the \emph{geometry of maps}---integral extensions and the Nullstellensatz (\S\S9--10)---and
the \emph{finiteness and decomposition} theory of Noetherian rings (\S11). The climax is
\emph{dimension} (\S12), refined locally by \emph{completion} (\S13). We cash out in
dimension one---\emph{DVRs and Dedekind domains}, the arithmetic payoff (\S14)---and close
by pointing to the \emph{horizons}: schemes, sheaves, and the Langlands programs (\S15).
\end{ppreview}

\begin{nnote}[Why a number theorist or a geometer should keep reading]
Every object in these notes is load-bearing in modern research. Completions give the
$p$-adic numbers $\ints_p$ and $\mathbb{Q}_p$; the \emph{local fields} of the
\vocab{local Langlands} correspondence are exactly these completions, and the \emph{adeles}
of \vocab{global Langlands} bundle all of them over $\Spec\ints$ at once. Dedekind domains
and their class groups are the start of class field theory; regular local rings and
dimension theory are the start of the geometry of singularities and \vocab{geometric
Langlands}. We will flag these connections as the relevant tools come online---the present
notes are, among other things, the commutative-algebra prerequisite for all of it.
\end{nnote}

\sources{A--M, Introduction and Ch.~1 (prime and maximal ideals, the nilradical, local
rings); A--K, Preface (the categorical/adjoint thesis) and \S\S1--3; Aluffi, the prefaces to
Ch.~III for the ``rings as the ground of geometry'' framing. Skim these for orientation, not
detail---the careful versions are \S\S3--4 below.}

% ======================================================================
\section{The categorical language, from scratch}
% ======================================================================

\noindent
The modern subject is written in the language of categories, so we learn just enough of it
right away---introduced gently, with rings and modules as the running examples, and assuming
no prior exposure. The payoff is enormous: a \vocab{universal property} lets us pin an
object down by how it maps to everything else (no formulas), and \vocab{adjoint functors}
organize half of what follows (tensor, localization, completion, free constructions) under
a single slogan. We are not doing category theory for its own sake; we are learning the
three or four ideas that the rest of the notes literally cannot be stated without.

\subsection{Categories, functors, and natural transformations}

We keep promising that a construction is ``natural,'' ``canonical,'' or that two objects are
``the same kind of thing.'' Category theory is the language that makes these words mean
something. A word of reassurance before we start: we are not studying category theory for
its own sake, and we need only its first few definitions. They earn their keep immediately,
because the central constructions of commutative algebra---tensor products, localization,
free modules---are cleanest, and are pinned down uniquely, by the categorical notions of
\emph{universal property} and \emph{adjoint functor} introduced over the next pages.

\begin{ddefinition}[Category]
A \vocab{category} $\cat{C}$ consists of: a collection of \vocab{objects}; for each pair of
objects $A,B$ a set of \vocab{morphisms} (or \vocab{arrows}) $\Hom_{\cat{C}}(A,B)$, written
$f\colon A\to B$; a \vocab{composition} law $\Hom(B,C)\times\Hom(A,B)\to\Hom(A,C)$,
$(g,f)\mapsto g\circ f$; and for each object $A$ an \vocab{identity} $\id_A\in\Hom(A,A)$.
These must satisfy two axioms: composition is associative, $h\circ(g\circ f)=(h\circ g)\circ
f$; and identities are neutral, $\id_B\circ f = f = f\circ\id_A$.
\end{ddefinition}

The definition is short because almost everything is an example. What it captures is
``objects and the allowed transformations between them, with a sensible notion of doing one
after another.''

\begin{eexample}[The categories we will live in]
\begin{itemize}\setlength{\itemsep}{2pt}
\item $\cat{Set}$: sets and functions.
\item $\cat{Grp}$, $\cat{Ab}$: groups (resp.\ abelian groups) and homomorphisms.
\item $\cat{Ring}$: commutative rings and ring homomorphisms (preserving $1$). \emph{This and
the next are the homes of this course.}
\item $R\text{-}\cat{Mod}$: modules over a fixed ring $R$ and $R$-linear maps. The central
category of these notes; when $R=k$ is a field this is $k$-vector spaces, and when
$R=\ints$ it is $\cat{Ab}$.
\item $\cat{Top}$: topological spaces and continuous maps---the target of $\Spec$.
\item \emph{A poset as a category.} Given a partially ordered set $(P,\le)$, let the objects
be the elements and put a single arrow $a\to b$ exactly when $a\le b$. Composition is
transitivity. For instance, the ideals of a ring $R$ ordered by inclusion form a category.
This example matters because it shows a ``morphism'' need not be a function at all---a
category is more flexible than ``sets with structure.''
\end{itemize}
\end{eexample}

\begin{iintuition}[It is all about the arrows]
The mental shift category theory asks for is this: stop looking \emph{inside} an object (its
elements, its construction) and look instead at the \emph{arrows} into and out of it. We
will see that an object is determined---up to unique isomorphism---by how it maps to and from
everything else. That is the engine behind ``universal properties'' below, and it is why we
can define the tensor product without ever choosing a basis.
\end{iintuition}

Categories are related by functors, which are ``maps of categories.''

\begin{ddefinition}[Functor]
A \vocab{functor} $F\colon\cat{C}\to\cat{D}$ assigns to each object $A$ of $\cat{C}$ an
object $F(A)$ of $\cat{D}$, and to each morphism $f\colon A\to B$ a morphism
$F(f)\colon F(A)\to F(B)$, preserving identities ($F(\id_A)=\id_{F(A)}$) and composition
($F(g\circ f)=F(g)\circ F(f)$). It is \vocab{covariant} as stated; a \vocab{contravariant}
functor instead reverses arrows, $F(f)\colon F(B)\to F(A)$, with $F(g\circ f)=F(f)\circ
F(g)$.\footnote{Equivalently, a contravariant functor on $\cat{C}$ is a covariant functor on
the \emph{opposite category} $\cat{C}^{\mathrm{op}}$, which has the same objects but all
arrows reversed.}
\end{ddefinition}

Functors are not exotic; the operations you already perform on rings and modules are mostly
functors, and recognizing them as such is what lets us reason about them uniformly.

\begin{eexample}[The operations of commutative algebra are functors]
\begin{itemize}\setlength{\itemsep}{2pt}
\item \emph{Forgetful functors}, e.g.\ $\cat{Ring}\to\cat{Set}$ (forget the ring structure)
or $R\text{-}\cat{Mod}\to\cat{Ab}$ (forget scalar multiplication).
\item \emph{Free functors}, e.g.\ $\cat{Set}\to R\text{-}\cat{Mod}$ sending a set $S$ to the
free module $R^{(S)}$. (Soon: left adjoint to the forgetful functor.)
\item \emph{The spectrum} $\Spec\colon\cat{Ring}\to\cat{Top}$, \emph{contravariant}---a ring
map $R\to S$ yields $\Spec S\to\Spec R$, as in \S1.
\item \emph{The $\Hom$ functors.} For a fixed $R$-module $N$, $\Hom_R(N,-)$ is a covariant
functor $R\text{-}\cat{Mod}\to R\text{-}\cat{Mod}$, while $\Hom_R(-,N)$ is contravariant.
\item \emph{Tensoring} $-\otimes_R N$ is a covariant functor (\S7), and \emph{localization}
$S^{-1}(-)$ is a covariant functor $R\text{-}\cat{Mod}\to S^{-1}R\text{-}\cat{Mod}$ (\S6).
\end{itemize}
The whole later drama---left versus right exactness, $\Tor$ and $\Ext$---is about how these
last functors fail to preserve exact sequences. That is why we set them up as functors now.
\end{eexample}

Finally, functors themselves can be compared, and the comparison is where the word
``natural'' acquires its technical meaning.

\begin{ddefinition}[Natural transformation]
Given two functors $F,G\colon\cat{C}\to\cat{D}$, a \vocab{natural transformation}
$\eta\colon F\Rightarrow G$ is a choice, for every object $A$, of a morphism
$\eta_A\colon F(A)\to G(A)$, such that for every morphism $f\colon A\to B$ the square
\[
\begin{array}{ccc}
F(A) & \xrightarrow{\ \eta_A\ } & G(A)\\[2pt]
{\scriptstyle F(f)}\big\downarrow & & \big\downarrow{\scriptstyle G(f)}\\[2pt]
F(B) & \xrightarrow{\ \eta_B\ } & G(B)
\end{array}
\]
commutes. If every $\eta_A$ is an isomorphism, $\eta$ is a \vocab{natural isomorphism} and we
write $F\cong G$.
\end{ddefinition}

\begin{iintuition}[``Natural'' means ``uniformly, without choices'']
A natural transformation is the same construction carried out at every object, compatibly
with all morphisms. This is the rigorous meaning of ``canonical.'' The textbook contrast is
from linear algebra: for finite-dimensional vector spaces the map $V\to V^{**}$ into the
double dual is natural (it needs no basis and commutes with every linear map), whereas the
map $V\to V^{*}$ to the dual is \emph{not} natural---you can only build one by choosing a
basis, and different choices give incompatible maps. ``Canonical'' was always trying to mean
``natural''; now it does. A purely algebraic instance, which you will meet in \S6: for a
multiplicative set $S\subseteq R$ the localization maps $M\to S^{-1}M$, $m\mapsto m/1$, are
natural in $M$---the same construction at every module, compatible with every $R$-linear
map---so localization comes packaged with a canonical natural transformation
$\id\Rightarrow S^{-1}(-)$ from the identity functor.
\end{iintuition}

\begin{xca}[\textbf{(easy)} $\Hom$ is a functor]
Fix an $R$-module $N$. Check that $\Hom_R(N,-)$ is a covariant functor: describe what it does
to a morphism $g\colon M\to M'$, and verify it preserves identities and composition. Then do
the same for the contravariant $\Hom_R(-,N)$.
\end{xca}

\begin{xca}[\textbf{(classic)} natural versus chosen]
On finite-dimensional $k$-vector spaces, show $V\mapsto V^{*}=\Hom_k(V,k)$ is a contravariant
functor and $V\mapsto V^{**}$ is a covariant functor, and that the evaluation maps
$\mathrm{ev}_V\colon V\to V^{**}$ assemble into a natural isomorphism. Explain in one
sentence why there is no natural isomorphism $V\cong V^{*}$. \emph{(This is the cleanest
possible illustration of the definition.)}
\end{xca}

\subsection{Universal properties: characterizing an object by its maps}

Here is the single most useful technique category theory gives us. A \vocab{universal
property} characterizes an object by the morphisms into or out of it---never by its internal
construction. The reason to care is twofold: a universal property defines a construction
\emph{without arbitrary choices}, and it automatically proves the construction is
\emph{unique up to a unique isomorphism}. After we state a universal property we will always
also give an explicit model (a concrete construction) and check the two describe the same
object; the universal property tells us \emph{what} the thing is, the model proves it
\emph{exists}.

The simplest universal properties name an object by a single mapping condition.

\begin{ddefinition}[Initial and terminal objects]
An object $I$ of $\cat{C}$ is \vocab{initial} if for every object $A$ there is exactly one
morphism $I\to A$. It is \vocab{terminal} (or \vocab{final}) if for every $A$ there is
exactly one morphism $A\to I$. An object that is both is a \vocab{zero object}.
\end{ddefinition}

\begin{eexample}[Initial and terminal, in our categories]
In $\cat{Set}$: the empty set $\varnothing$ is initial, any one-point set is terminal. In
$R\text{-}\cat{Mod}$: the zero module is both---a zero object. In $\cat{Ring}$: the
\emph{zero ring} is terminal (every ring maps onto it), and---crucially---$\ints$ is
\emph{initial}, because every ring $R$ receives exactly one ring homomorphism from $\ints$
(it must send $1\mapsto 1$, hence $n\mapsto n\cdot 1_R$, and this is forced). We return to
how special this makes $\ints$ in \S2.3.
\end{eexample}

The reason universal properties are trustworthy is the following two-line theorem, which we
prove because it is the engine that makes ``\emph{the} tensor product,'' ``\emph{the}
localization,'' etc.\ well-defined.

\begin{ttheorem}[Universal properties pin objects down]
Any two initial objects of a category are isomorphic, by a \emph{unique} isomorphism. (Dually
for terminal objects.)
\end{ttheorem}

\begin{proof}
Let $I,I'$ be initial. Initiality of $I$ gives a unique map $u\colon I\to I'$; initiality of
$I'$ gives a unique $v\colon I'\to I$. Then $v\circ u\colon I\to I$ is a map from $I$ to
itself, but initiality of $I$ says there is only \emph{one} such map, and $\id_I$ is one, so
$v\circ u=\id_I$. Symmetrically $u\circ v=\id_{I'}$. Hence $u$ is an isomorphism, and it was
the unique map $I\to I'$ to begin with.
\end{proof}

\begin{ttechnique}[How to read and use a universal property]
Every universal property in this book is secretly the statement ``\emph{this object is
initial (or terminal) in an auxiliary category of solutions to a mapping problem}.'' So it
always buys you the same three things: (i) a recipe---``a map out of $X$ is the same as the
following data''; (ii) uniqueness up to unique isomorphism, for free, by the theorem above;
and (iii) a way to \emph{construct} maps out of (or into) $X$ without touching its internals.
Don't ask what the object \emph{is}; ask what maps to and from it \emph{do}.
\end{ttechnique}

We now meet the workhorse universal properties. In each, notice the rhythm: \emph{universal
property}, then \emph{concrete model}, then \emph{why they agree}.

\begin{eexample}[Quotient ring $R/I$: the universal way to kill $I$]
\emph{Universal property.} The quotient map $\pi\colon R\to R/I$ is the universal ring
homomorphism out of $R$ that sends $I$ to $0$: for every ring map $\varphi\colon R\to S$ with
$\varphi(I)=0$ there is a \emph{unique} $\bar\varphi\colon R/I\to S$ with
$\bar\varphi\circ\pi=\varphi$. \emph{Concrete model.} Cosets $r+I$, with the inherited
operations. \emph{Why they agree.} The map $\bar\varphi(r+I)\defeq\varphi(r)$ is
well-defined exactly because $\varphi(I)=0$, and it is the only possible value---this is the
first isomorphism theorem, read as a universal property.\footnote{Aluffi, Ch.~III \S3,
develops quotient rings precisely this way; the same pattern (a quotient as a universal
``map that kills a subobject'') recurs for modules in \S4 below.}
\end{eexample}

\begin{eexample}[Free module $R^{(S)}$: maps out of it are just functions on $S$]
\emph{Universal property.} For a set $S$ there is a module $R^{(S)}$ with an inclusion
$S\hookrightarrow R^{(S)}$ such that every set-function $S\to M$ into an $R$-module $M$
extends \emph{uniquely} to an $R$-linear map $R^{(S)}\to M$. \emph{Concrete model.} Finite
formal combinations $\sum_{s} r_s\, e_s$ ($r_s\in R$, almost all $0$). \emph{Why they agree.}
``Define a linear map by saying where the generators go, then extend linearly'' is exactly
the unique extension---this is the basis-free phrasing of the most familiar move in linear
algebra.
\end{eexample}

\begin{eexample}[{Polynomial ring $R[x]$: ``plugging in'' is a universal property}]
\emph{Universal property.} $R[x]$ is the free commutative $R$-algebra on one generator: for
any $R$-algebra $S$, an $R$-algebra map $R[x]\to S$ is the \emph{same data} as a choice of
one element of $S$ (the image of $x$). \emph{Concrete model.} Polynomials $\sum a_i x^i$.
\emph{Why they agree.} Sending $x\mapsto s$ and extending by the algebra laws is the
substitution (``evaluation'') homomorphism, and it is the unique extension---this is precisely
\emph{why} you are allowed to plug anything you like into a polynomial.\footnote{Aluffi,
Ch.~III \S2.2, ``Universal property of polynomial rings.'' This example is the prototype for
the \emph{functor of points} we mention in the note below.}
\end{eexample}

\begin{nnote}[A first glimpse of the functor of points]
The polynomial example says $\Hom_{\cat{Ring}}(R[x],S)\cong S$ \emph{naturally in $S$}: the
ring $R[x]$ ``is'' the rule that assigns to each ring $S$ its underlying set of elements---the
\emph{affine line}. This is Grothendieck's \vocab{functor of points} idea: an affine scheme
$\Spec R$ is encoded by the functor $S\mapsto\Hom_{\cat{Ring}}(R,S)$ of its ``$S$-valued
points.'' We are not using this yet, but it is the reason the categorical language is not
optional in modern algebraic geometry (\S15).
\end{nnote}

\begin{xca}[\textbf{(easy)} the two ends of $\cat{Ring}$]
Prove carefully that $\ints$ is initial and the zero ring is terminal in $\cat{Ring}$. Where
exactly does ``ring homomorphisms preserve $1$'' get used?
\end{xca}

\begin{xca}[\textbf{(classic)} the quotient as a universal property]
State the universal property of $R/I$ precisely and prove it from the coset construction.
Then deduce the first isomorphism theorem (any $\varphi\colon R\to S$ factors as
$R\twoheadrightarrow R/\ker\varphi\xrightarrow{\sim}\im\varphi\hookrightarrow S$) from it.
\end{xca}

\subsection{Limits, colimits, and the initial/terminal objects (why \texorpdfstring{$\ints$}{Z} is special)}

Initial and terminal objects are the smallest universal properties; products, quotients,
intersections, and unions are all instances of two general patterns called \vocab{limits} and
\vocab{colimits}. We introduce them lightly---enough to recognize the constructions of
commutative algebra as members of these two families, because that recognition is what makes
the adjoint-functor theorems of \S2.4 so powerful.

\begin{ddefinition}[Limit and colimit, informally]
A \vocab{diagram} in $\cat{C}$ is a collection of objects and morphisms (a picture of arrows).
A \vocab{limit} of the diagram is a universal object mapping \emph{to} all of it: an object
$L$ with compatible maps to every object of the diagram, through which any other such
``compatible family of maps'' factors uniquely. A \vocab{colimit} is the dual: a universal
object receiving maps \emph{from} all of it. (Limits generalize ``intersection/product''; colimits
generalize ``union/glue.'')
\end{ddefinition}

\begin{eexample}[The (co)limits you already use]
\begin{itemize}\setlength{\itemsep}{2pt}
\item \emph{Products and coproducts.} The \vocab{product} $\prod A_i$ is the limit of a
diagram with no arrows (it has projections to each $A_i$); the \vocab{coproduct} is the
colimit (it has inclusions from each $A_i$). In $R\text{-}\cat{Mod}$ a finite product is the
direct sum $\bigoplus A_i$, which is \emph{also} the coproduct---a biproduct. In $\cat{Ring}$
the product is the usual $R\times S$, while the coproduct is the tensor product
$R\otimes_\ints S$ (\S7).
\item \emph{Kernels and cokernels.} In $R\text{-}\cat{Mod}$, $\ker f$ is a limit and
$\coker f$ is a colimit; quotient modules are colimits.
\item \emph{Direct limits} $\dirlim$, colimits over a directed system, are pervasive in
commutative algebra: every module is the direct limit of its finitely generated submodules,
and (\S6) the localization $R_\mfp=\dirlim_{f\notin\mfp}R_f$ is a direct limit.
\item \emph{Inverse limits} $\invlim$ are the limits we complete with: the $p$-adic integers
are $\ints_p=\invlim_n \ints/p^n\ints$, and more generally $\widehat{R}=\invlim_n R/I^n$
(\S13).
\end{itemize}
\end{eexample}

\noindent
Not every category has all such (co)limits on hand---the category of fields, for instance, has
no products, since a product $k\times k'$ of two fields is not a field (it has zero-divisors).
What makes $R\text{-}\cat{Mod}$ and $\cat{Ring}$ so pleasant to work in is precisely that they
\emph{are} complete and cocomplete: every diagram in them has both a limit and a colimit. We
will use this silently and often.

Now we can say precisely why $\ints$ deserved a starring role in \S1, and why number theory
is, structurally, ``geometry over a single base point.''

\begin{iidea}[Why $\ints$ is special, and why everything lives over $\Spec\ints$]
$\ints$ is the \emph{initial} object of $\cat{Ring}$: every ring receives a unique ring map
from $\ints$. Translating through the contravariant dictionary of \S1, this says
$\Spec\ints$ is the \emph{terminal} object among affine schemes---\emph{every} commutative
ring, hence every affine scheme, has a unique map to $\Spec\ints$. So all of commutative
algebra and algebraic geometry takes place ``over $\Spec\ints$,'' and the special case of
rings that happen to be algebras over a field $k$ is the case of geometry over the point
$\Spec k$. Number theory is then exactly the geometry of $\Spec\ints$ and its covers (\S1's
$\ints[i]$ example); this ``arithmetic = geometry over the terminal base $\Spec\ints$'' is
the conceptual seed of Arakelov theory and much of modern arithmetic geometry.
\end{iidea}

\begin{xca}[\textbf{(easy)} products and coproducts by hand]
Identify the product and the coproduct of two objects in each of $\cat{Set}$, $\cat{Ab}$, and
$\cat{Ring}$. (You should find disjoint union vs.\ cartesian product in $\cat{Set}$; direct
sum $=$ direct product for two objects in $\cat{Ab}$; and $R\times S$ vs.\
$R\otimes_\ints S$ in $\cat{Ring}$.)
\end{xca}

\begin{xca}[\textbf{(medium)} completion and localization as (co)limits]
Exhibit the structure maps making $\ints_p=\invlim_n\ints/p^n\ints$ an inverse limit, and
(looking ahead to \S6) describe the directed system whose direct limit is $R_\mfp$. In each
case, what is the universal property doing for you?
\end{xca}

\subsection{Adjoint functors: the single most important idea in the book}

Everything so far converges here. \vocab{Adjoint functors} are the organizing principle of
the subject: tensor products, localization, free modules, base change, and completion are all
one half of an adjunction, and two short theorems about adjunctions will let us predict their
behavior toward exact sequences before we compute anything. We motivate the definition with
the example you have already seen twice.

Recall the free module: a set-function $S\to U(M)$ (where $U$ forgets the module structure)
is the same data as an $R$-linear map $R^{(S)}\to M$. Spelled out, there is a bijection,
natural in both $S$ and $M$,
\[
\Hom_{R\text{-}\cat{Mod}}\!\big(R^{(S)},\,M\big)\;\cong\;\Hom_{\cat{Set}}\!\big(S,\,U(M)\big).
\]
This bijection \emph{is} an adjunction between the free functor $F\colon S\mapsto R^{(S)}$ and
the forgetful functor $U$.

\begin{ddefinition}[Adjoint functors]
Functors $F\colon\cat{C}\to\cat{D}$ and $G\colon\cat{D}\to\cat{C}$ form an \vocab{adjunction},
with $F$ \vocab{left adjoint} to $G$ (and $G$ \vocab{right adjoint} to $F$), written
$F\dashv G$, if there is a bijection
\[
\Hom_{\cat{D}}\!\big(F(c),\,d\big)\;\cong\;\Hom_{\cat{C}}\!\big(c,\,G(d)\big)
\]
natural in $c\in\cat{C}$ and $d\in\cat{D}$.
\end{ddefinition}

\begin{iintuition}[Left adjoints build, right adjoints forget]
Picture $G$ as ``forgetting'' some structure and $F$ as ``freely building it back.'' The
adjunction says: to map a freely-built object $F(c)$ into $d$ is no more and no less than to
map the raw material $c$ into the forgotten shadow $G(d)$---you specify the map on generators
and the structure forces the rest. Left adjoints are the \emph{constructive, generators-and-
relations} functors (free modules, tensor, localization, base change); right adjoints are the
\emph{forgetful, limit-like} functors (the forgetful functor, $\Hom$). Keeping this picture
straight tells you which side of every construction you are on.
\end{iintuition}

The reason this single definition reorganizes the subject is the following list, every entry
of which we will use.

\begin{eexample}[The adjunctions that run commutative algebra]
\begin{itemize}\setlength{\itemsep}{3pt}
\item \emph{Free $\dashv$ forgetful} (the model): $R^{(-)}\dashv U$ for modules, and
$\ints[-]\dashv U$ for rings (the polynomial ring $\ints[S]$ is the free commutative ring on a
set $S$).
\item \emph{Tensor $\dashv$ Hom} (``adjoint associativity''), the most-used adjunction in the
book:
\[
\Hom_R\!\big(M\otimes_R N,\,P\big)\;\cong\;\Hom_R\!\big(M,\,\Hom_R(N,P)\big),
\]
i.e.\ $-\otimes_R N\ \dashv\ \Hom_R(N,-)$ (\S7).
\item \emph{Base change $\dashv$ restriction of scalars}: for a ring map $R\to S$, the
functor $S\otimes_R-$ is left adjoint to restricting an $S$-module to an $R$-module. (\S7.)
\item \emph{Localization $\dashv$ restriction}: $S^{-1}(-)$ is left adjoint to restriction
along $R\to S^{-1}R$---localization is the special case of base change with $S^{-1}R$ (\S6).
\end{itemize}
\end{eexample}

What makes these worth naming all at once is that adjointness, by itself, predicts how each
functor treats limits and colimits---and hence exact sequences.

\begin{ttheorem}[The slogan that organizes \S\S5--8]
Left adjoints preserve colimits (in particular cokernels, direct sums, and direct
limits);\footnote{The proof is a one-line calculation in the spirit of the ``it is all about
the arrows'' box. If $F\dashv G$ and $c=\dirlim c_i$, then for every object $d$ we have
$\Hom(F(\dirlim c_i),d)\cong\Hom(\dirlim c_i, Gd)\cong\invlim\Hom(c_i,Gd)\cong
\invlim\Hom(Fc_i,d)\cong\Hom(\dirlim Fc_i,d)$, naturally in $d$ (the middle step is that
$\Hom(-,Gd)$ turns colimits into limits). Since an object is pinned down by the maps out of
it---the \vocab{Yoneda lemma}---this forces $F(\dirlim c_i)\cong\dirlim Fc_i$. The statement
for right adjoints is the dual.}
right adjoints preserve limits (in particular kernels, products, and inverse limits).
Consequently $-\otimes_R N$ is right exact and $\Hom_R(N,-)$, $\Hom_R(-,N)$ are left
exact.\footnote{This is exactly the structural observation A--K put at the center of their
treatment: from ``Tensor is the left adjoint of Hom'' and ``a left adjoint preserves direct
limits/colimits'' they deduce right-exactness of tensor and much else (A--K, Preface and
\S\S6--9). A--M prove the same isomorphism (their ``adjoint associativity'') but, writing
before this language was standard, do not name it an adjunction; we do.} Moreover adjoints,
when they exist, are unique up to natural isomorphism---so ``\emph{the}'' tensor product and
``\emph{the}'' localization are well-defined.
\end{ttheorem}

\begin{eequation}[Carry this for the next four sections]
\[
\textbf{left adjoint}\ \Rightarrow\ \text{right exact (preserves colimits)},
\]
\[
\textbf{right adjoint}\ \Rightarrow\ \text{left exact (preserves limits)}.
\]
Tensor and localization are left adjoints; $\Hom$ is a right adjoint. The \emph{failure} of
the missing half---how far $\otimes$ is from left exact, how far $\Hom$ is from right
exact---is measured by the derived functors $\Tor$ and $\Ext$ (\S8). Half of homological
algebra is bookkeeping for this one slogan.
\end{eequation}

\begin{nnote}[The same idea, all the way up to Langlands]
Adjunctions are not special to commutative algebra; they are how mathematics expresses
``free/forgetful'' and ``induction/restriction'' everywhere. In representation theory,
induction and restriction of representations are adjoint (Frobenius reciprocity); in the
\vocab{Langlands} program their analogues---parabolic induction and Jacquet
restriction---are adjoint functors, and the geometric Langlands correspondence is, at bottom,
an equivalence between two (derived) categories implemented by functors of exactly this kind.
The language you are learning here is the same language those theories are written in.
\end{nnote}

\begin{xca}[\textbf{(easy)} the free module adjunction]
Write out the natural bijection
\[
\Hom_{R\text{-}\cat{Mod}}\big(R^{(S)},M\big)\;\cong\;\Hom_{\cat{Set}}\big(S,U(M)\big)
\]
explicitly in both directions, and check it is natural in $M$. Conclude $R^{(-)}\dashv U$.
\end{xca}

\begin{xca}[\textbf{(classic)} adjoint associativity]
Construct natural maps in both directions between $\Hom_R(M\otimes_R N,P)$ and
$\Hom_R(M,\Hom_R(N,P))$ and show they are mutually inverse. \emph{(You may take the
generators-and-relations construction of $\otimes$ as known from \S7; the point is to see the
adjunction concretely.)}
\end{xca}

\begin{xca}[\textbf{(medium)} right-exactness for free]
Granting the theorem that left adjoints preserve cokernels, deduce that $-\otimes_R N$ is
right exact: applying it to $M'\to M\to M''\to 0$ gives an exact
$M'\otimes N\to M\otimes N\to M''\otimes N\to 0$. Where, precisely, did ``left adjoint'' do
the work?
\end{xca}

\begin{xca}[\textbf{(thinker)} adjoints are essentially unique]
Show that if $F$ has two right adjoints $G,G'$ then $G\cong G'$ naturally. \emph{(Hint: both
represent the same functor $d\mapsto\Hom_{\cat{D}}(F(-),d)$; use the uniqueness theorem of
\S2.2.)} Conclude that ``the'' tensor product and ``the'' localization are well-defined.
\end{xca}

\sources{Aluffi Ch.~I \S\S3--5 (categories, morphisms, universal properties; initial/final
objects, products, coproducts) is the gentle, example-rich source for \S\S2.1--2.3; Ch.~VIII
\S1 covers functors and natural transformations, and \S2 the tensor/$\Tor$ adjunction. A--K
threads adjoints through the whole book---its Preface and \S\S6--9 are the place to watch
``left adjoint $\Rightarrow$ right exact'' do real work. For category theory in its own
right, S.\ Mac Lane, \emph{Categories for the Working Mathematician}, is the standard
reference.}

% ======================================================================
\section{Rings, ideals, and the spectrum}
% ======================================================================

\noindent
Now the geometry begins. We review rings, ideals, and quotients quickly---but with a modern
slant---and then build the central object: the \vocab{prime spectrum} $\Spec R$, topologized
by the \vocab{Zariski topology}. The conceptual leap, and the one worth slowing down for,
is \emph{why the points are the prime ideals} (not just the maximal ones) and how the
elements of $R$ act as functions on $\Spec R$. We make $\Spec$ into a contravariant functor
so that algebra and geometry translate automatically, and we draw the first pictures.

\subsection{Rings, homomorphisms, ideals, and quotients: a fast modern review}

\subsection{Prime and maximal ideals, and why the primes are the points}

\subsection{The Zariski spectrum and its topology}

\subsection{Nilradical, Jacobson radical, and the radical of an ideal}

\subsection{\texorpdfstring{$\Spec$}{Spec} as a contravariant functor: ring maps are space maps}

\subsection{First pictures: the spectra of \texorpdfstring{$\ints$, $k[x]$, $k[x,y]$}{Z, k[x], k[x,y]}, and a fat point}

\sources{A--M Ch.~1; A--K \S\S1--3 (Rings, Prime Ideals, Radicals); Aluffi Ch.~III \S\S1--4.}

% ======================================================================
\section{Modules: linear algebra over a ring}
% ======================================================================

\noindent
A ring is best understood through the things it acts on, and those are its \vocab{modules}.
Modules are ``vector spaces over a ring,'' and the slogan of the modern subject is that
commutative algebra is the study of the category $R\text{-}\mathbf{Mod}$. Almost everything
we already know is a module: ideals, quotients, abelian groups (= $\ints$-modules), vector
spaces (= modules over a field). The one new phenomenon---and the source of the whole
subject's richness---is that \emph{not every module is free}; measuring that failure is
what occupies us for the rest of the notes. We introduce annihilators and support to begin
seeing modules as objects spread out over $\Spec R$.

\subsection{Modules and the category \texorpdfstring{$R\text{-}\mathbf{Mod}$}{R-Mod}: why everything is a module}

\subsection{Submodules, quotients, \texorpdfstring{$\Hom$}{Hom}, and the operations on submodules}

\subsection{Generators and relations; finitely generated and finitely presented modules}

\subsection{Free modules, and the productive failure of ``every module is free''}

\subsection{Annihilators and support: a first glimpse of modules as sheaves on \texorpdfstring{$\Spec R$}{Spec R}}

\sources{A--M Ch.~2; A--K \S4 (Modules); Aluffi Ch.~III \S5 and Ch.~VI.}

% ======================================================================
\section{Exact sequences and the homological reflex}
% ======================================================================

\noindent
Exact sequences are how modern algebra bookkeeps relationships between modules, and learning
to read them fluently is half the battle. We cover short and long exact sequences, the
\vocab{snake lemma} (the workhorse of every diagram chase), and---crucially---the idea that
a functor's \emph{failure} to preserve exactness is information, not a nuisance. That single
reframing is what turns $\Hom$ and $\otimes$ into the derived functors $\Ext$ and $\Tor$
later, and it is the reflex that distinguishes the modern treatment from the classical one.

\subsection{Kernels, cokernels, images; short and long exact sequences}

\subsection{The snake lemma and the art of the diagram chase}

\subsection{Extensions and splitting: what an exact sequence remembers}

\subsection{Additive, left-exact, and right-exact functors as a measuring stick}

\sources{A--M Ch.~2 (exact sequences); A--K \S5; Aluffi Ch.~VI \S7 and Ch.~VIII--IX.}

% ======================================================================
\section{Localization: zooming in on the spectrum}
% ======================================================================

\noindent
Localization is the most important \emph{technique} in the subject, and intuition is the
whole battle, so it gets room to breathe. Geometrically, inverting a set of elements means
\emph{zooming in} on part of $\Spec R$; in particular, localizing at a prime $\mfp$ produces
the \vocab{local ring} $R_\mfp$, which records exactly the behavior of $R$ near the point
$\mfp$. We give both the universal property and the explicit fractions construction, say why
they describe the same object, and then state the \vocab{local--global principle}: a great
many properties hold for $R$ if and only if they hold at every prime. That principle is the
engine that lets us prove global facts one point at a time.

\subsection{The problem: studying a ring near a point, or inverting a set of elements}

\subsection{The construction \texorpdfstring{$S^{-1}R$ and $S^{-1}M$}{S\^{}-1 R}, and its universal property}

\subsection{The two cases that matter: a basic open \texorpdfstring{$R_f$}{R\_f} and a local ring \texorpdfstring{$R_\mfp$}{R\_p}}

\subsection{Local rings, residue fields, and why ``local'' is the right word}

\subsection{Localization is exact; the order-preserving correspondence of primes}

\subsection{The local--global principle: prove it at every prime}

\sources{A--M Ch.~3 (this is the heart of the book---read it twice); A--K \S\S11--13
(Localization of Rings, of Modules, and Support).}

% ======================================================================
\section{Tensor products, base change, and flatness}
% ======================================================================

\noindent
The tensor product is where the categorical language earns its keep. We give the universal
property (the clean definition) \emph{and} the generators-and-relations construction (the
one you compute with), and we say precisely why they are the same object. The key structural
fact is \vocab{adjoint associativity}: $-\otimes_R N$ is left adjoint to $\Hom_R(N,-)$, from
which right-exactness falls out for free. \vocab{Flatness}---the modules for which tensoring
preserves exactness---is the algebraic shadow of ``continuously varying family,'' and it is
one of the deepest geometric ideas in the subject. We close with Lazard's theorem, the
modern characterization of flat as a filtered colimit of finite free modules.

\subsection{Bilinear maps, the universal property, and the explicit construction (why they agree)}

\subsection{Tensor as a functor; right-exactness from adjoint associativity}

\subsection{Extension of scalars and base change; localization \emph{is} base change}

\subsection{Flatness: the modules that preserve exactness, i.e.\ ``continuous families''}

\subsection{Lazard's theorem: flat \texorpdfstring{$=$}{=} filtered colimit of finite free modules}

\sources{A--M Ch.~2--3; A--K \S\S6--9 (Direct Limits, Filtered Direct Limits, Tensor
Products, Flatness)---the most thoroughly modern treatment of the three; Aluffi Ch.~VIII \S2.}

% ======================================================================
\section{Derived functors: Tor and Ext}
% ======================================================================

\noindent
Once we accept that $\otimes$ and $\Hom$ fail to be exact, the natural question is: by how
much, and what does the failure measure? The answer is \vocab{homological algebra}. We build
projective and injective resolutions, define \vocab{$\Tor$} (the derived functor of $\otimes$)
and \vocab{$\Ext$} (the derived functor of $\Hom$), and read their meaning: $\Tor$ detects
flatness and controls how subvarieties intersect, while $\Ext^1$ classifies extensions and
is the home of obstructions. This is the toolkit behind depth, regular sequences, and the
homological characterization of smoothness that reappears in the dimension theory below.

\subsection{Projective and injective modules, and resolutions}

\subsection{\texorpdfstring{$\Tor$}{Tor}: the homology of \texorpdfstring{$\otimes$}{tensor}, and what its vanishing means}

\subsection{\texorpdfstring{$\Ext$}{Ext}: extensions, obstructions, and \texorpdfstring{$\Ext^1$}{Ext\^{}1} as a classifying group}

\subsection{Homological algebra as the modern measure of ``how far from exact''}

\sources{A--K (the flatness and Tor threads in \S9 and beyond); Aluffi Ch.~VIII \S\S2--6 and
Ch.~IX (Homological algebra). For depth and the Auslander--Buchsbaum circle, see Eisenbud
later.}

% ======================================================================
\section{Integral extensions and the geometry of finite maps}
% ======================================================================

\noindent
Integral extensions are the algebraic side of \emph{finite maps} of spaces---think of a
branched cover of curves. We develop integral elements, integral closure, and \vocab{normal}
rings, then the three transfer theorems---\vocab{lying over}, \vocab{going up}, and
\vocab{going down}---which say how primes (points) move under a finite map. \vocab{Noether
normalization}, the statement that every finitely generated algebra over a field is finite
over a polynomial subring, is both a structural cornerstone and the lever we use to prove the
Nullstellensatz in the next section.

\subsection{Integral elements, integral extensions, integral closure, and normal rings}

\subsection{Lying over, going up, going down: how points move under a finite map}

\subsection{Noether normalization: every affine algebra is finite over a polynomial ring}

\subsection{Geometric meaning: finite surjections, branching, and normalization}

\sources{A--M Ch.~5; A--K \S10 (Cayley--Hamilton---the slick route to integrality), \S14
(Cohen--Seidenberg, i.e.\ going up/down), \S15 (Noether Normalization).}

% ======================================================================
\section{The Nullstellensatz: the dictionary made precise}
% ======================================================================

\noindent
Here the algebra--geometry dictionary becomes a theorem. \vocab{Hilbert's
Nullstellensatz} identifies the maximal ideals of $k[x_1,\dots,x_n]$ (for $k$ algebraically
closed) with the points of affine space, and upgrades the loose correspondence ``ideals
$\leftrightarrow$ vanishing sets'' into the precise statement $I(V(\mfa))=\sqrt{\mfa}$. The
upshot is an honest equivalence between affine varieties and reduced finitely generated
$k$-algebras---the prototype of the ring/space duality and the doorway to scheme theory.

\subsection{Hilbert's Nullstellensatz, weak and strong}

\subsection{The radical--variety correspondence \texorpdfstring{$\sqrt{\mfa}=I(V(\mfa))$}{sqrt(a)=I(V(a))}}

\subsection{Coordinate rings and the equivalence varieties \texorpdfstring{$\simeq$}{\textasciitilde} reduced f.g.\ algebras}

\subsection{Why this is the gateway from varieties to schemes}

\sources{A--M Ch.~5 and Ch.~7; A--K \S15 (Noether Normalization gives the cleanest proof);
field-theory background in Aluffi Ch.~VII.}

% ======================================================================
\section{Chain conditions: Noetherian rings and primary decomposition}
% ======================================================================

\noindent
Almost every ring that arises in geometry or number theory satisfies the \vocab{ascending
chain condition}---it is \vocab{Noetherian}---and \vocab{Hilbert's basis theorem} guarantees
this is preserved by the constructions we care about. The structural payoff is
\vocab{primary decomposition}, the ring-theoretic analogue of factoring an integer, which we
present \emph{the modern way}: \vocab{associated primes} of a module from the very start
(not, as in older books, extracted afterward from a decomposition). The associated primes,
together with the support, encode the irreducible components and the embedded points of a
scheme---geometry read straight off the algebra.

\subsection{Noetherian and Artinian modules; the ascending chain condition}

\subsection{Hilbert's basis theorem: why everything in sight is Noetherian}

\subsection{Associated primes, defined from the start (the modern definition)}

\subsection{Primary decomposition of modules and the uniqueness theorems}

\subsection{\texorpdfstring{$\Ass$}{Ass}, \texorpdfstring{$\Supp$}{Supp}, and the geometry of components and embedded points}

\sources{A--M Ch.~4 (primary decomposition), Ch.~6--7 (chain conditions, Noetherian rings);
A--K \S\S16--18 (Chain Conditions, Associated Primes, Primary Decomposition)---note A--K puts
associated primes \emph{first}, as is now standard.}

% ======================================================================
\section{Krull dimension: counting the directions a space can go}
% ======================================================================

\noindent
How many independent directions does a space have? The algebraic answer is \vocab{Krull
dimension}: the longest chain of prime ideals. The deep and satisfying theorem is that three
\emph{a priori} different notions of dimension---chains of primes, the growth rate of the
\vocab{Hilbert function}, and the minimal number of parameters cutting out a point---all
agree for Noetherian local rings. \vocab{Krull's principal ideal theorem} pins down
codimension one, and \vocab{regular local rings} give the first clean algebraic criterion for
\emph{smoothness} versus singularity.

\subsection{Krull dimension via chains of primes; height and codimension}

\subsection{Hilbert functions, Hilbert polynomials, and dimension by growth}

\subsection{Systems of parameters and the dimension theorem (the three notions agree)}

\subsection{Krull's principal ideal theorem and codimension one}

\subsection{Regular local rings: a first algebraic criterion for smoothness}

\sources{A--M Ch.~11 (Dimension Theory); A--K \S\S19--21 (Length, Hilbert Functions,
Dimension).}

% ======================================================================
\section{Completion: the local-analytic picture}
% ======================================================================

\noindent
Localization zooms in to a point; \vocab{completion} goes one step further and records the
\emph{infinitesimal, analytic} neighborhood of that point. Built as an inverse limit
$\widehat{R}=\invlim R/I^{n}$, completion produces objects like the $p$-adic integers
$\ints_p$ and the formal power series $k[[x]]$, where convergence is free and
\vocab{Hensel's lemma} lets us lift approximate solutions to exact ones. Completion is a flat
base change, so it imports the homological machinery intact, and the associated graded ring
links it back to the Hilbert-function dimension theory just developed.

\subsection{Filtrations, the \texorpdfstring{$I$}{I}-adic topology, and inverse limits}

\subsection{Examples: \texorpdfstring{$\ints_p$, $k[[x]]$}{Z\_p, k[[x]]}, and formal neighborhoods}

\subsection{Hensel's lemma and lifting solutions; why complete local rings are ``analytic''}

\subsection{Completion as a flat base change; the associated graded ring}

\sources{A--M Ch.~10 (Completions; graded rings and the Artin--Rees lemma); A--K \S22.}

% ======================================================================
\section{Dimension one: DVRs, Dedekind domains, fractional ideals}
% ======================================================================

\noindent
The one-dimensional case is where commutative algebra pays its great debt to number theory.
A \vocab{discrete valuation ring} is the algebra of a single smooth point or a single prime,
and a \vocab{Dedekind domain} is a ring---like the ring of integers of a number field---in
which \emph{ideals} factor uniquely even when \emph{elements} do not. \vocab{Fractional
ideals} and the \vocab{ideal class group} measure precisely how badly unique factorization
fails, turning the classical headache of $\ints[\sqrt{-5}]$ into a clean invariant. This is
the concrete arithmetic payoff of the whole machine.

\subsection{Discrete valuation rings: the algebra of one smooth prime}

\subsection{Dedekind domains: unique factorization of ideals, not elements}

\subsection{Fractional ideals and the ideal class group}

\subsection{Rings of integers, \texorpdfstring{$\ints[\sqrt{-5}]$}{Z[sqrt(-5)]}, and the arithmetic payoff}

\sources{A--M Ch.~9 (DVRs and Dedekind Domains); A--K \S\S23--26 (Discrete Valuation Rings,
Dedekind Domains, Fractional Ideals, Arbitrary Valuation Rings).}

% ======================================================================
\section{Where this goes: schemes, sheaves, and the Langlands programs}
% ======================================================================

\noindent
A short outlook on what the path opens onto, with no claim to completeness. Gluing the affine
pieces $\Spec R$ along their localizations gives \vocab{schemes}; promoting modules to
$\Spec R$ globally gives \vocab{quasi-coherent sheaves} and their cohomology---this is the
launchpad for algebraic geometry proper. On the arithmetic side, viewing $\Spec\ints$ as a
``curve'' and a ring of integers as a branched cover of it is the geometric heart of class
field theory and, beyond it, the global, local, and geometric \vocab{Langlands programs},
where every tool we built (localization, completion, Dedekind domains, $\Spec$) is load-bearing.

\subsection{From \texorpdfstring{$\Spec$}{Spec} to schemes: gluing affines and the structure sheaf}

\subsection{Quasi-coherent sheaves: modules globalized, and their cohomology}

\subsection{The arithmetic horizon: number fields and \texorpdfstring{$\Spec\ints$}{Spec Z} as a curve}

\subsection{A word on the Langlands programs (global, local, geometric) and where CA sits}

\sources{Outlook only---see the ``what to read next'' list: Hartshorne and Vakil for schemes,
Neukirch for the arithmetic, and the survey literature for Langlands.}

% ======================================================================
\section{End-of-path exercises}
% ======================================================================

\noindent
A short, path-spanning set, meant to be attempted after a full read of the notes built on
this skeleton. Difficulty tags as above; the ``thinkers'' deliberately push toward the
references.

\begin{xca}[\textbf{(easy)} reading the dictionary]
For $R=k[x,y]/(xy)$, describe $\Spec R$ as a topological space, list its minimal primes, and
explain---using only the radical---why $R$ is reduced but not a domain. Then say, in one
sentence each, what ``reduced'' and ``not a domain'' mean \emph{geometrically}.
\end{xca}

\begin{xca}[\textbf{(classic)} localization is base change]
Let $S\subseteq R$ be multiplicative and $M$ an $R$-module. Construct a natural isomorphism
$S^{-1}M \cong S^{-1}R \otimes_R M$ and use it, together with right-exactness of $\otimes$, to
re-prove that localization is exact. Where, exactly, does flatness of $S^{-1}R$ enter?
\end{xca}

\begin{xca}[\textbf{(medium)} the local--global reflex]
Show that an $R$-module homomorphism $f\colon M\to N$ is injective (resp.\ surjective, resp.\
an isomorphism) if and only if $f_\mfp\colon M_\mfp\to N_\mfp$ is for every prime $\mfp$. State
carefully which direction needs $M$ or $N$ finitely generated, and give a one-line geometric
gloss of what the statement says about sections over $\Spec R$.
\end{xca}

\begin{xca}[\textbf{(medium)} dimension three ways, on one example]
For the local ring $R=k[x,y]_{(x,y)}$, compute the Krull dimension by chains of primes, by
the degree of the Hilbert--Samuel polynomial, and by exhibiting a system of parameters. Verify
all three give $2$, and identify the maximal ideal's minimal number of generators to check
that $R$ is regular.
\end{xca}

\begin{xca}[\textbf{(thinker)} where unique factorization goes]
In $R=\ints[\sqrt{-5}]$, exhibit the failure of unique factorization on $6$, then \emph{repair}
it: factor the ideal $(6)$ into prime ideals and identify the relevant ideal classes. Explain
how the class group converts a defect of \emph{elements} into an invariant of the \emph{ring},
and point to where Dedekind domains made this work. \emph{This is the gateway to algebraic
number theory; aim to internalize the mechanism, not just the computation.}
\end{xca}

% ======================================================================
\section*{References and a reading plan}
\addcontentsline{toc}{section}{References and a reading plan}
% ======================================================================

\noindent
These notes are written against three primary sources, listed here in the order I would reach
for them. Together they \emph{are} a reading plan: A--M for the lean classical spine, A--K for
the same material redone in fully modern categorical language, and Aluffi for the gentle,
example-rich on-ramp to the categorical and module-theoretic ideas.
\sloppy

\medskip
\noindent\textbf{Primary sources (the three this path is built on).}
\begin{itemize}\setlength{\itemsep}{3pt}
\item \textbf{M.\ F.\ Atiyah \& I.\ G.\ Macdonald}, \emph{Introduction to Commutative
Algebra}. The standard concise reference; the entire subject in $130$ pages, with the famous
exercises. Terse but never wrong---read it with pen in hand and fill the gaps yourself. Best
for the lean logical spine (\S\S3--14 of this path).
\item \textbf{A.\ Altman \& S.\ Kleiman}, \emph{A Term of Commutative Algebra} (freely
available). A--M rewritten for the modern era: associated primes from the start, direct limits
and adjoints front and center, every result tied to a universal property, and complete
solutions to every exercise. The best single match for the ``modern, categorical'' framing of
these notes (especially \S\S2, 6--8).
\item \textbf{P.\ Aluffi}, \emph{Algebra: Chapter 0}. Categorical and unusually readable; the
right place to learn universal properties, the module/linear-algebra perspective, and
$\Tor$/$\Ext$ without prior homological algebra. Best for \S\S2, 4--5, 8.
\end{itemize}

\medskip
\noindent\textbf{What to read next / books to own.}
\begin{itemize}\setlength{\itemsep}{3pt}
\item \textbf{D.\ Eisenbud}, \emph{Commutative Algebra with a View Toward Algebraic Geometry}.
The big modern book: everything here, plus homological depth, regular sequences, and the
geometry made explicit. The one to grow into.
\item \textbf{H.\ Matsumura}, \emph{Commutative Ring Theory}. The standard advanced reference
for dimension theory, completion, and excellent rings.
\item \textbf{M.\ Reid}, \emph{Undergraduate Commutative Algebra}. A short, friendly,
geometry-forward companion to read in parallel with A--M.
\item \textbf{R.\ Hartshorne}, \emph{Algebraic Geometry} (Ch.~II) and \textbf{R.\ Vakil},
\emph{The Rising Sea: Foundations of Algebraic Geometry} (\url{math.stanford.edu/~vakil/216blog/}).
Where \S15 actually goes: schemes and quasi-coherent sheaves.
\item \textbf{J.\ Neukirch}, \emph{Algebraic Number Theory}. The arithmetic continuation of
\S14---Dedekind domains, ramification, and $\Spec\ints$ as a curve.
\item \textbf{The Stacks Project} (\url{stacks.math.columbia.edu}). A complete, hyperlinked,
ground-up development of commutative algebra and algebraic geometry; the modern reference of
record.
\end{itemize}

\end{document}
